\begin{explanationbox}
\textbf{\textcolor{CExplanation}{先看图形：三条切线围出什么？}}

想象两根竖直的栏杆，分别固定在 $x=-a$ 和 $x=a$；椭圆上的动点 $P$ 处再放一根恰好贴住椭圆的切线。它与左右栏杆交于 $P_1,P_2$。从长轴端点 $A_1,A_2$ 向上（或向下）到交点的长度分别是 $h_1,h_2$。

切线倾斜时，一边可能变长，另一边就会变短；它们不是\textbf{和不变}，而是\textbf{积不变}：$h_1h_2=b^2$。所以一个变成 $2b$，另一个必变成 $b/2$。

\medskip
\textbf{\textcolor{CExplanation}{画图时最容易看错的地方}}
\begin{itemize}[leftmargin=*,itemsep=0.4em,topsep=0.3em]
\item $A_1P_1$、$A_2P_2$ \textbf{沿竖直方向}，二者是梯形的两条平行底边；二底之间的水平距离恒为 $2a$。
\item $P_1,P_2$ 永远在 $x$ 轴的\textbf{同一侧}；当 $P$ 在椭圆上半部，它们都在上方；在下半部，它们都在下方。
\item 两边一样高时最省面积：$h_1=h_2=b$，图形成为宽 $2a$、高 $b$ 的矩形，面积 $2ab$。注意\textbf{并非椭圆的面积 $\pi ab$}。
\end{itemize}

\medskip
\textbf{\textcolor{CExplanation}{直觉一：为什么一长一短，而且恰好互为倒数？}}

设 $P=(a\cos\theta,b\sin\theta)$，其中 $\sin\theta\ne0$。椭圆切线化成
\[
\frac{x\cos\theta}{a}+\frac{y\sin\theta}{b}=1.
\]
分别令 $x=-a$ 和 $x=a$，得
\[
y_1=b\frac{1+\cos\theta}{\sin\theta}
=b\cot\frac{\theta}{2},\qquad
y_2=b\frac{1-\cos\theta}{\sin\theta}
=b\tan\frac{\theta}{2}.
\]
原来两个长度就是同一角的 $\tan$ 与 $\cot$ 乘上 $b$！而 $\tan u\cdot\cot u=1$，所以 $|y_1y_2|=b^2$。这说明所谓“神奇定值”，本质上就是\textbf{互为倒数的两个量}。

\medskip
\textbf{\textcolor{CExplanation}{直觉二：把椭圆想成被横向拉长的圆}}

先画一个半径为 $b$ 的圆 $u^2+v^2=b^2$，再把所有点的横坐标放大 $a/b$ 倍：
\[
(x,y)=\left(\frac ab u,v\right).
\]
圆变成所给椭圆；圆在 $u=\pm b$ 的两条竖直切线变成椭圆在 $x=\pm a$ 的切线。\textbf{竖直长度完全没变}，于是圆上的“竖直长度乘积为 $b^2$”也被原封不动地保留下来。面积则横向放大 $a/b$ 倍：圆的最小面积 $2b^2$ 变成 $2ab$。

注意：仿射拉伸保留直线、交点、切线关系及平行性，\textbf{一般不保留斜线长度和角度}；所以只能把圆上的竖直距离直接搬过来，不能把“圆的等长切线段”不加证明地搬到椭圆上。

\medskip
\textbf{\textcolor{CExplanation}{直觉三：为什么面积最小值几乎不用计算？}}

面积只看 $h_1+h_2$，定值只看 $h_1h_2=b^2$。回想基本不等式
\[
h_1+h_2\ge2\sqrt{h_1h_2}=2b.
\]
又因为两条竖直底边间的距离是 $2a$，面积就是 $a$ 倍的高度和：
\[
S=a(h_1+h_2)\ge2ab.
\]
\textbf{最小的和来自两个相等的正数。}当 $h_1=h_2$ 时动点切线水平，椭圆上仅在上下两个短轴端点具有水平切线。这就解释了取等位置。

\medskip
\textbf{\textcolor{CExplanation}{一个坐标图像：斜率控制不平衡，截距控制面积}}

把切线记成 $y=kx+m$：
\[
y_1=m-ak,\qquad y_2=m+ak.
\]
两侧高度以 $m$ 为中心，上下各偏离 $ak$；其积是平方差 $m^2-a^2k^2$。椭圆恰好要求切线满足
\[
m^2=a^2k^2+b^2.
\]
因此\textbf{高度积固定，截距绝对值 $|m|$ 决定面积}：
\[
S=2a|m|\ge2ab.
\]
可以记成“\textbf{端点代入 $\pm a$，乘积凑平方差；面积取两边平均高}”。

\medskip
\textbf{\textcolor{CExplanation}{与熟悉结论的联系}}
\begin{itemize}[leftmargin=*,itemsep=0.4em,topsep=0.3em]
\item \textbf{圆的切线：} 圆上同一点的切线构造更容易用“同一外点的两条切线段等长”来证；椭圆可由圆水平拉伸得到。
\item \textbf{切线判别式：} 直线 $y=kx+m$ 与椭圆相切 $\Longleftrightarrow m^2=a^2k^2+b^2$，能最快做选填题。
\item \textbf{均值不等式：} 两个正长度\textbf{积固定时，和最小等于两倍几何平均数}。本题是这一模式的标准几何化。
\item \textbf{对称变式：} 把竖直端点切线改成水平端点切线，交换 $x,y$、$a,b$，定值由 $b^2$ 变为 $a^2$。
\end{itemize}

\medskip
\textbf{\textcolor{CExplanation}{考场识别与记忆}}

题中同时出现\textbf{“椭圆两轴端点处的平行切线”与“第三条动切线”}，就优先想：两个截距相乘、定值，再利用基本不等式求和的最小值。
\[
\boxed{\text{端点 }\pm a\ \to\ \text{平方差}\ \to\ b^2
\quad;\quad \text{积固定}\ \to\ \text{和最小}.}
\]
\end{explanationbox}
